% answers_short.tex —— 2019 年·天津卷·理科 · 简短答案速查
% 由 tools/make_short_answers.py 自 answers.tex 的 \gksolhead/\gkanswer 提取，
% 并与 2019/2019_数学真题_天津卷理_答案要点.md 逐题比对后生成，勿手改（改源 MD 后重跑）。

\begin{gknotice}
\textbf{使用说明}：本节只给最终答案便于快速核对，完整推导见解析版。答案为编写者独立推导，\textbf{非官方标准答案}。
\end{gknotice}

\gkpart{选择题}
\begin{center}\begin{tabular}{ccccccccc}
\toprule
题号 & 1 & 2 & 3 & 4 & 5 & 6 & 7 & 8 \\
\midrule
答案 & D & C & B & B & D & A & C & C \\
\bottomrule
\end{tabular}\end{center}

\gkpart{填空题}
\begin{itemize}[leftmargin=2.6em]
  \item[\textbf{9.}] $2 - 3\mathrm{i}$
  \item[\textbf{10.}] $28$
  \item[\textbf{11.}] $\dfrac{\pi}{4}$
  \item[\textbf{12.}] $\dfrac{3}{4}$
  \item[\textbf{13.}] $4\sqrt{3}$
  \item[\textbf{14.}] $-1$
\end{itemize}

\gkpart{解答题}
\begin{itemize}[leftmargin=2.6em]
  \item[\textbf{15.}] (1) $-\dfrac{1}{4}$；(2) $-\dfrac{3\sqrt{5} + 7}{16}$
  \item[\textbf{16.}] (1) 分布列见解析版，$E(X) = 2$；(2) $\dfrac{20}{243}$
  \item[\textbf{17.}] (1) 证明见解析版；(2) $\dfrac{4}{9}$；(3) $\dfrac{8}{7}$
  \item[\textbf{18.}] (1) $\dfrac{x^2}{5} + \dfrac{y^2}{4} = 1$；(2) $\pm\dfrac{2\sqrt{30}}{5}$
  \item[\textbf{19.}] (1) $a_n = 3n + 1$，$b_n = 3 \cdot 2^n$；(2) ① $a_{2^n}(c_{2^n} - 1) = 9 \cdot 4^n - 1$；② $\sum_{i=1}^{2^n} a_i c_i = \dfrac{27}{2} \cdot 4^n + 5 \cdot 2^{n-1} - n - 12$
  \item[\textbf{20.}] (1) 单调递增区间为 $\left[2k\pi - \dfrac{3\pi}{4}, 2k\pi + \dfrac{\pi}{4}\right]$ ($k \in \mathbb{Z}$)，单调递减区间为 $\left[2k\pi + \dfrac{\pi}{4}, 2k\pi + \dfrac{5\pi}{4}\right]$ ($k \in \mathbb{Z}$)；(2) 证明见解析版；(3) 证明见解析版
\end{itemize}
